\section{Protection closure and median feeds}\label{sec:paper-closure}
The incoming-core reduction takes induced subgraphs. We now keep the vertex
set fixed and add arcs. Every convenience predicate in this section refers
to the graph present at the stated stage.

\subsection{Two safe operations}
A missing direction $a\to b$ is \emph{forced} in $H$ if some vertex $w$
satisfies $a\triangleleft_H w$ and $w\to b\in A(H)$.

\begin{lemma}\label{lem:paper-forced-step}
Forced directions are consistent, uniquely convenient in $H$, and mandatory
in every tournament completion preserving all protections of $H$.
Adding all forced directions simultaneously preserves every old protection.
\end{lemma}
\begin{proof}
For a witness $a\triangleleft_H w$, $w\to b$, every predecessor $p$ of $a$
also points to $w$, and hence reaches $b$ within two steps. Thus $a\to b$
is convenient. Its reverse fails at $w$, since $w$ cannot reach its protected
far head $a$. Opposite directions cannot both be forced: each would have
been proved convenient while the other witness proves it nonconvenient.
In a tournament, $a\triangleleft w$ also gives outgoing domination, so
$w\to b$ forces $a\to b$.

Fix an old $y\triangleleft_H x$. An old predecessor of $y$ already points
to $x$. For a newly incoming $a\to y$, choose its witness
$a\triangleleft_H w$, $w\to y$. Then $w\to x$ originally. If $ax$ is
missing it is forced $a\to x$ in the same round. If $x\to a$ were an old
arc, protection at $a,w$ would give $x\to w$, a digon with $w\to x$.
Thus the required new incoming nesting is preserved in every case.
\end{proof}

There is a second, sequential operation. If $ab$ is missing,
$N_H^-(a)\subseteq N_H^-(b)$, and $b$ has no protection successor in $H$,
add the single arc $a\to b$. Call this a \emph{weak step}.

\begin{lemma}\label{lem:paper-weak-step}
A weak step is convenient in its preceding graph, preserves all old
protections, and creates $a\triangleleft b$.
\end{lemma}
\begin{proof}
Every old predecessor of $a$ already points to $b$. Only $N^-(b)$ changes.
An old relation $y\triangleleft b$ becomes easier to satisfy, and an old
relation $b\triangleleft x$ is excluded by the hypothesis on $b$.
All other incoming comparisons are unchanged. The added arc together with
the old containment gives $a\triangleleft b$.
\end{proof}

Starting at $D$, repeatedly perform a full forced round; when no forced
direction remains, perform one available weak step and recompute.
Stop when neither operation is available, obtaining a partial graph $Q$.
Any deterministic choice among eligible weak steps may be used.
The process terminates after at most $M(D)$ new pairs. Protections acquired
at any stage persist. In particular, every added arc has a tail with a
protection successor in $Q$. Simultaneous arbitrary weak choices are not
part of the construction.

\subsection{A fixed partial graph and a complete reference family}
Since $Q$ has no remaining forced pair, each $y\triangleleft_Q x$ has
outgoing as well as incoming nesting. Indeed, $x\to z$ excludes $z\to y$
by incoming nesting; if $yz$ were missing, $y\to z$ would still be forced.
A two-position swap therefore shows that every global median order of $Q$
extends protection: reversing $x,y$ gains their mutual arc and loses no
contribution from an intermediate vertex. Missing arcs contribute zero.

Complete every remaining missing pair forward in a linear extension of
the protection order of $Q$. Every such tournament preserves protection.
To see this, consider $y\triangleleft_Q x$ and a newly incoming $a\to y$.
The reference has $a<y<x$. An old $x\to a$ would have forced $y\to a$,
contrary to terminality and the missing pair $ay$. Thus $a\to x$ is either
old or is also completed forward. This proves the incoming nesting.

Let $T_*$ maximize the ordinary median score over the family of
\emph{all} these extensions. A remaining pair of $Q$ is called free.

\begin{theorem}[Rigidity and original neighbourhoods]\label{thm:paper-feed-invariance}
Every global median of $T_*$ respects protection in $Q$ and puts every free
arc forward. At any feed $f$ of such an order,
\[
 N_{T_*}^+(f)=N_Q^+(f)=N_D^+(f),\qquad
 N_Q^{++}(f)=N_D^{++}(f).
\]
If $r$ is the number of free pairs, the maximum tournament median score is
$\med(Q)+r$, and the possible feeds across the maximizing completions are
exactly the median feeds of $Q$.
\end{theorem}
\begin{proof}
Tournament domination makes every median respect protection. If a free arc
were backward in one median of $T_*$, use that median as a new reference.
This fixes the arcs of $Q$ and turns all backward free arcs forward,
strictly increasing the score. This contradicts maximality and applies to
every median of the same maximizing tournament. The argument requires the
entire extension family.

A feed has no outgoing free arc, since it would be backward. Nor can it be
the tail of an arc added during closure: its persistent protection successor
would have to occur later. Thus its first set was unchanged at every stage.
Induct over the stages. Any new two-step path at $f$ uses an original first
arc and a newly added second arc. Convenience in the preceding graph says
that this original predecessor $f$ already reached the head within two steps.
A first neighbour stays first throughout, so no new exact second head occurs.
Old exact second paths persist because the first set is fixed.

A median of $Q$ is an admissible reference and scores $\med(Q)+r$ after
completion. No order in any reference tournament can exceed this number:
its $Q$ contribution is at most $\med(Q)$ and its free contribution at most
$r$. Equality forces both bounds to be attained, proving the score and feed
assertions.
\end{proof}

\subsection{Equality sedimentation and the original far budget}
For a median order $L$ of a tournament with feed $f$, a vertex
$y\in N^-(f)$ is good if some $a\in N^+(f)$ occurs before $y$ in $L$ and
has $a\to y$. Let $G_L$ be the good set. The strong feed theorem of
Havet and Thomass\'e gives
\[
 |N^+(f)|\le |G_L|\le |N^{++}(f)|.
\]
When $|G_L|=|N^+(f)|$, their global sedimentation lemma replaces $L$ by
the order consisting of its bad vertices, then $f$, then $N^+(f)\cup G_L$,
preserving each block's previous relative order. It remains a global median
of the same tournament~\cite{HavetThomasse2000}. Each vertex of the last
block moves strictly later. These statements concern ordinary global medians.

\begin{theorem}[Strict original surplus]\label{thm:paper-strict-surplus}
Let $f$ be a median feed of $T_*$ with $g_D(f)>0$, and set
\[
 A_f=N_{T_*}^{++}(f)\setminus N_D^{++}(f),\qquad
 R_f=\mathcal M_Q(f),\quad m_f=|R_f|.
\]
Then $m_f>0$ and $t_D(f)\ge g_D(f)+1$.
If $g_{T_*}(f)=0$, then
\[
 \varnothing\ne R_f\subseteq U_{T_*}(f)\cap
       \bigl(\mathcal F_D(f)\setminus A_f\bigr),\qquad
 t_D(f)\ge g_D(f)+m_f.
\]
No all-deficiency or minimality assumption on the other vertices is needed.
\end{theorem}
\begin{proof}
Original protection survives. By Theorem~\ref{thm:paper-feed-invariance},
\[
 N_{T_*}^{++}(f)=N_D^{++}(f)\,\dot\cup\,A_f,
 \quad A_f\subseteq\mathcal F_D(f),\quad
 g_{T_*}(f)=g_D(f)-|A_f|\le0.
\]
If $R_f$ were empty, $f$ would be whole in $Q$. Every far head of a whole
vertex is protected in that graph. Preservation of these protections would
prevent any further second head at $f$, giving $g_{T_*}(f)=g_D(f)>0$.
This contradiction proves $m_f>0$.

If the completed gap is zero, strong feed equality gives
$G_L=N_{T_*}^{++}(f)$. A partner $r\in R_f\cap G_L$ has a free incoming
arc $r\to f$, since the first set at $f$ is original. Sedimentation puts
$f$ before $r$ while keeping a median, contradicting free-arc rigidity.
Thus every remaining partner is far in $T_*$. It was originally missing,
so it cannot be an originally protected head, and old exact second paths
survive. Hence it lies in $\mathcal F_D(f)\setminus A_f$. The two disjoint
sets have sizes $m_f$ and $g_D(f)$. If the completed gap is negative instead,
integrality gives $|A_f|\ge g_D(f)+1$ directly.
\end{proof}

\begin{theorem}[An original partner cycle]\label{thm:paper-original-cycle}
Under the zero-gap hypotheses of Theorem~\ref{thm:paper-strict-surplus},
every vertex of $Q[R_f]$ has an incoming neighbour there, and the original
graph $D[R_f]$ contains a directed cycle. Consequently
\[
 m_f\ge3,\qquad t_D(f)\ge g_D(f)+m_f\ge g_D(f)+3.
\]
\end{theorem}
\begin{proof}
The feed is maximal in terminal protection. For $r\in R_f$, terminality
of the weak operation gives $p\in N_Q^-(r)\setminus N_Q^-(f)$.
Since $r$ is far from $f$ in $T_*$, $p\to r$ implies $p\to f$ there.
Neither direction on $pf$ was present in $Q$, so $p\in R_f$. Thus each
remaining partner has a predecessor in the induced terminal graph.

Fix the final set $R_f$ throughout the construction. For a forced internal
addition $a\to b$, let $a\triangleleft_H w$, $w\to b$ be its preceding
witness. Far-ness of $b$ in $T_*$ gives $w\to f$ there. If $wf$ were adjacent
in $Q$, its direction would already be $w\to f$, and the retained protection
$a\triangleleft w$ would force the still-missing $a\to f$. Therefore
$w\in R_f$. The internal addition shortcuts an old internal two-step path.
All simultaneous such shortcuts preserve acyclicity of an acyclic induced
graph.

For a single weak internal addition $a\to b$, creating the first cycle
would require an old internal path $b\leadsto p\to a$. It cannot have length
one, because $ab$ was missing. Old incoming containment gives $p\to b$,
already completing an old cycle. Thus weak steps cannot create the first
cycle either. An acyclic $D[R_f]$ would remain acyclic, whereas a finite
graph of positive minimum indegree has a cycle. Orientedness excludes cycle
lengths one and two.
\end{proof}

The stronger $+3$ bound applies to zero completed gap only. Omitting all weak
steps still gives Theorem~\ref{thm:paper-strict-surplus}, but does not give
the partner-cycle conclusion. The original cycle in the theorem is not a
claim that every terminal cycle lifts edge by edge, nor that its vertices
are actual bad-direction sources.

\begin{corollary}\label{cor:paper-low-far}
Every oriented graph with $t_D(v)\le1$ at every vertex has a Seymour vertex.
More generally this holds if $t_D(v)\le g_D(v)$ at every deficient vertex.
\end{corollary}
\begin{proof}
Otherwise every original gap is positive and any feed contradicts the
strict-surplus theorem.
\end{proof}

If $h_D(f)$ counts the protected original far heads, partitioning the other
vertices gives $t_D(f)-g_D(f)=n-1-2d_D^+(f)-h_D(f)$. Hence every
original-positive-gap closure feed satisfies
\begin{equation}\label{eq:paper-degree-localization}
 2d_D^+(f)+h_D(f)\le n-2.
\end{equation}
At $n=2\delta^+(D)+3$, such a feed has minimum original outdegree and at
most one protected far head. This locates a possibly deficient feed; it
does not assert that a minimum-degree vertex already has the SNP in $D$.

\subsection{The negative-gap equality branch}
Every new original-far head belongs to $G_L$: its second arc cannot lie in
$Q$, and therefore is free and forward, putting its original first intermediate
before the head. Write
$b_L=|N_D^{++}(f)\setminus G_L|$ and
$s_L=|G_L|-d_D^+(f)$. Then
\begin{equation}\label{eq:paper-good-slack}
 s_L=|A_f|-g_D(f)-b_L\ge0.
\end{equation}
When $s_L=0$, sedimentation excludes $R_f\cap G_L$. Its second partners
number at most $|A_f|-g_D(f)$ and are original nongood second vertices;
its far partners number at most $t_D(f)-|A_f|$. Thus
\begin{equation}\label{eq:paper-equality-partners}
 |R_f|\le t_D(f)-g_D(f)\qquad(s_L=0).
\end{equation}
This includes negative completed gap. Put
$R_2=R_f\cap N_D^{++}(f)$, $R_\infty=R_f\cap U_{T_*}(f)$, and
$U=R_\infty\setminus\Reach_{D[R_f]}(R_2)$, where reachability includes
the seed vertices. The supplement proves that closure creates no new
reachability from $R_2$ into $R_\infty$ inside $R_f$. If the whole set $U$
is nonempty, its original induced graph $D[U]$ contains a directed cycle
(Theorem~\ref{S-thm:weak-equality-original-partner-path-cycle}).
When $s_L>0$, the cited equality-sedimentation move is unavailable. This
does not assert that every median order of the tournament has the same slack.
